\documentclass[11pt]{article}
\usepackage[margin=1in]{geometry}
\usepackage{amsmath, amssymb}
\usepackage{enumitem}
\usepackage{booktabs}
\usepackage{array}
\usepackage{fancyhdr}
\usepackage{hyperref}

\pagestyle{fancy}
\fancyhf{}
\lhead{CEC 315 --- Signals and Systems}
\rhead{Homework 6 --- Official Solutions}
\cfoot{\thepage}

\title{\textbf{CEC 315 Homework 6 --- Official Solutions} \\
  \large The $z$-Transform: Lectures 19--21 \\
  \normalsize Spring 2026 \quad Instructor: Rogelio Gracia Otalvaro}
\author{}
\date{}

% Helper macros
\newcommand{\zt}{\xleftrightarrow{\mathcal{Z}}}
\newcommand{\ztu}{\xleftrightarrow{\mathcal{Z}_u}}
\newcommand{\ROC}[1]{\text{ROC: } #1}

\begin{document}
\maketitle
\thispagestyle{fancy}

\noindent\textit{This document is the official instructor solution set, faithfully transcribed from
\texttt{hw\_practice\_problems/hw-lctr19-21-solutions.pdf} (12 pages). It is provided
alongside the earlier student-made file \texttt{hw6\_solutions.tex} for side-by-side comparison.
Final answers are boxed; regions of convergence (ROCs) are explicitly stated for every
$z$-transform.}

\bigskip

\noindent\textbf{Key $z$-transform pairs used throughout:}
\begin{align*}
a^{n}u[n] &\;\zt\; \frac{1}{1-az^{-1}}, & \ROC{|z|>|a|}, \\
-a^{n}u[-n-1] &\;\zt\; \frac{1}{1-az^{-1}}, & \ROC{|z|<|a|}, \\
n\,a^{n}u[n] &\;\zt\; \frac{az^{-1}}{(1-az^{-1})^{2}}, & \ROC{|z|>|a|}, \\
\delta[n-k] &\;\zt\; z^{-k}, & \\
r^{n}\cos(\omega_0 n)\,u[n] &\;\zt\; \frac{1-r\cos\omega_0\,z^{-1}}{1-2r\cos\omega_0\,z^{-1}+r^{2}z^{-2}}, & \ROC{|z|>r}.
\end{align*}

%=============================================================
\section*{Problem 1: Computing $z$-Transforms and ROCs}
%=============================================================

\subsection*{Part (a): $x_1[n] = 5(0.7)^{n}u[n]$}
Using $a^{n}u[n]\zt \dfrac{1}{1-az^{-1}}$ with $a=0.7$:
\begin{equation*}
\boxed{X_1(z)=\dfrac{5}{1-0.7\,z^{-1}}, \qquad |z|>0.7.}
\end{equation*}
Pole at $z=0.7$. Check: $X_1(1)=5/0.3=50/3$.

\subsection*{Part (b): $x_2[n] = -(4)^{n}u[-n-1]$}
Using $-a^{n}u[-n-1]\zt \dfrac{1}{1-az^{-1}}$ with $a=4$, ROC $|z|<|a|$:
\begin{equation*}
\boxed{X_2(z)=\dfrac{1}{1-4\,z^{-1}}, \qquad |z|<4.}
\end{equation*}
Pole at $z=4$. Check: $X_2(1)=1/(1-4)=-1/3$.

Direct: $\displaystyle\sum_{n=-\infty}^{-1}-(4)^{n}=-\sum_{m=1}^{\infty}4^{-m}=-\frac{1/4}{1-1/4}=-\frac{1}{3}.\;\checkmark$

\subsection*{Part (c): $x_3[n] = 3\delta[n]-2\delta[n-1]+\delta[n-4]$}
Using $\delta[n-k]\zt z^{-k}$:
\begin{equation*}
\boxed{X_3(z)=3-2z^{-1}+z^{-4}, \qquad |z|>0.}
\end{equation*}
Finite-duration sequence; ROC is all $z\neq 0$. Check: $X_3(1)=3-2+1=2$.

\subsection*{Part (d): $x_4[n] = n(0.5)^{n}u[n]$}
Using $n\,a^{n}u[n]\zt \dfrac{az^{-1}}{(1-az^{-1})^{2}}$ with $a=0.5$:
\begin{equation*}
\boxed{X_4(z)=\dfrac{0.5\,z^{-1}}{(1-0.5\,z^{-1})^{2}}, \qquad |z|>0.5.}
\end{equation*}
Double pole at $z=0.5$, zero at $z=0$. Check: $X_4(1)=0.5/0.25=2$.

\subsection*{Part (e): $x_5[n] = 2(0.3)^{n}u[n]+4(0.9)^{n}u[n]$}
\textbf{Step 1 --- transform each term:}
\[
\frac{2}{1-0.3\,z^{-1}},\;|z|>0.3 \qquad\text{and}\qquad \frac{4}{1-0.9\,z^{-1}},\;|z|>0.9.
\]
\textbf{Step 2 --- common denominator:}
\[
X_5(z)=\frac{2(1-0.9\,z^{-1})+4(1-0.3\,z^{-1})}{(1-0.3\,z^{-1})(1-0.9\,z^{-1})}=\frac{6-3.0\,z^{-1}}{(1-0.3\,z^{-1})(1-0.9\,z^{-1})}.
\]
ROC: $|z|>0.3 \cap |z|>0.9 = |z|>0.9$.
\begin{equation*}
\boxed{X_5(z)=\dfrac{6-3\,z^{-1}}{(1-0.3\,z^{-1})(1-0.9\,z^{-1})}, \qquad |z|>0.9.}
\end{equation*}
Poles: $z=0.3,\;z=0.9$. Zero: $6-3z^{-1}=0\Rightarrow z=0.5$.

Check: $X_5(1)=\dfrac{6-3}{(0.7)(0.1)}=\dfrac{3}{0.07}=\dfrac{300}{7}\approx 42.86$. Direct: $\dfrac{2}{0.7}+\dfrac{4}{0.1}=\dfrac{300}{7}.\;\checkmark$

\subsection*{Part (f): $x_6[n] = (0.6)^{n}u[n]-3(2)^{n}u[-n-1]$}
\textbf{Step 1.} $(0.6)^{n}u[n]\zt\dfrac{1}{1-0.6\,z^{-1}}$, $|z|>0.6$.

For the second term, $-3(2)^{n}u[-n-1]=3\bigl[-(2)^{n}u[-n-1]\bigr]\zt\dfrac{3}{1-2\,z^{-1}}$, $|z|<2$.

\textbf{Step 2.}
\[
X_6(z)=\frac{1}{1-0.6\,z^{-1}}+\frac{3}{1-2\,z^{-1}}.
\]
ROC: $|z|>0.6\cap |z|<2=\{z:0.6<|z|<2\}$ (annulus).
\begin{equation*}
\boxed{X_6(z)=\dfrac{1}{1-0.6\,z^{-1}}+\dfrac{3}{1-2\,z^{-1}}, \qquad 0.6<|z|<2.}
\end{equation*}
Since the unit circle $|z|=1$ satisfies $0.6<1<2$, the unit circle lies in the ROC and the \textbf{DTFT exists}.$\;\checkmark$

\subsection*{Part (g): $x_7[n] = (0.8)^{n}\cos(0.25\pi n)\,u[n]$}
With $r=0.8$, $\omega_0=0.25\pi$: $r\cos\omega_0\approx 0.5657$, $2r\cos\omega_0\approx 1.1314$, $r^{2}=0.64$.
\begin{equation*}
\boxed{X_7(z)=\dfrac{1-0.5657\,z^{-1}}{1-1.1314\,z^{-1}+0.64\,z^{-2}}, \qquad |z|>0.8.}
\end{equation*}
Poles (complex conjugate pair): $z=0.8\,e^{\pm j 0.25\pi}$, magnitude $r=0.8$, angles $\pm 45^{\circ}$. Oscillation frequency $\omega_0=0.25\pi$ rad/sample. All poles have $|z|=0.8<1$, so the unit circle is in the ROC $|z|>0.8$; \textbf{DTFT exists}.$\;\checkmark$

\subsection*{Part (h): Sanity checks at $z=1$}
All checks were performed inline above. The rule $X(1)=\sum_n x[n]$ holds because $z^{-n}|_{z=1}=1$ for all $n$.

\begin{center}
\begin{tabular}{@{}clll@{}}
\toprule
Signal & $X(1)$ & $\sum x[n]$ & match \\
\midrule
(a) & $50/3\approx 16.67$ & $50/3$ & $\checkmark$ \\
(b) & $-1/3$ & $-1/3$ & $\checkmark$ \\
(c) & $2$ & $2$ & $\checkmark$ \\
(d) & $2$ & $2$ & $\checkmark$ \\
(e) & $300/7$ & $300/7$ & $\checkmark$ \\
(f) & $2.5-3=-0.5$ & $-0.5$ & $\checkmark$ \\
(g) & $\approx 0.854$ & $\approx 0.854$ & $\checkmark$ \\
\bottomrule
\end{tabular}
\end{center}

%=============================================================
\section*{Problem 2: Inverse $z$-Transform via Partial Fractions}
%=============================================================

\subsection*{Part (a): Distinct real poles, causal}
Given $X(z)=\dfrac{1+3z^{-1}}{(1-0.2z^{-1})(1-0.6z^{-1})}$, ROC $|z|>0.6$.

\textbf{Step 1 --- partial fractions:} $X(z)=\dfrac{A}{1-0.2z^{-1}}+\dfrac{B}{1-0.6z^{-1}}$.
\begin{itemize}[leftmargin=*]
\item $z^{-1}=5$: $1+15=A(1-3)\Rightarrow A=-8$.
\item $z^{-1}=5/3$: $1+5=B(1-1/3)=\tfrac{2}{3}B\Rightarrow B=9$.
\end{itemize}

\textbf{Step 2 --- directions from ROC.} ROC $|z|>0.6$ is outside both pole circles, so both terms are right-sided.

\textbf{Step 3:}
\begin{equation*}
\boxed{x[n]=\bigl[-8\,(0.2)^{n}+9\,(0.6)^{n}\bigr]\,u[n].}
\end{equation*}
Check: $x[0]=-8+9=1$. Initial-value theorem: $\lim_{z\to\infty}X(z)=1.\;\checkmark$

\subsection*{Part (b): Distinct real poles, two-sided signal}
Given $X(z)=\dfrac{4}{(1-0.5z^{-1})(1-3z^{-1})}$, ROC $0.5<|z|<3$.

\textbf{Step 1 --- partial fractions:}
\begin{itemize}[leftmargin=*]
\item $z^{-1}=2$: $4=A(1-6)=-5A\Rightarrow A=-4/5$.
\item $z^{-1}=1/3$: $4=B(1-1/6)=(5/6)B\Rightarrow B=24/5$.
\end{itemize}

\textbf{Step 2 --- assign directions from the annular ROC $0.5<|z|<3$:}
\begin{itemize}[leftmargin=*]
\item Pole at $z=0.5$: ROC is \textbf{outside} $\Rightarrow$ right-sided.
\item Pole at $z=3$: ROC is \textbf{inside} $\Rightarrow$ left-sided.
\end{itemize}

\textbf{Step 3 --- invert:}
\[
\tfrac{-4/5}{1-0.5z^{-1}}\;\longrightarrow\; -\tfrac{4}{5}(0.5)^{n}u[n],\qquad
\tfrac{24/5}{1-3z^{-1}}\;\longrightarrow\; -\tfrac{24}{5}(3)^{n}u[-n-1].
\]
\begin{equation*}
\boxed{x[n]=-\dfrac{4}{5}\,(0.5)^{n}u[n]-\dfrac{24}{5}\,(3)^{n}u[-n-1].}
\end{equation*}

\noindent\textbf{Key Point.} The left-sided inversion uses $\dfrac{B}{1-az^{-1}}$ with ROC $|z|<|a|$ $\longleftrightarrow$ $-B\,a^{n}u[-n-1]$. The negative sign is part of the pair itself.

\subsection*{Part (c): Repeated poles}
Given $X(z)=\dfrac{z^{-1}}{(1-0.5z^{-1})^{2}}$, ROC $|z|>0.5$.
\[
X(z)=\frac{z^{-1}}{(1-0.5z^{-1})^{2}}=\frac{1}{0.5}\cdot\frac{0.5\,z^{-1}}{(1-0.5z^{-1})^{2}}=2\cdot\frac{0.5\,z^{-1}}{(1-0.5z^{-1})^{2}}.
\]
\begin{equation*}
\boxed{x[n]=2n\,(0.5)^{n}u[n].}
\end{equation*}
Check: $x[0]=0,\; x[1]=1,\; x[2]=1$. Initial value $\lim_{z\to\infty}X(z)=0$, consistent.$\;\checkmark$

\subsection*{Part (d): Complex conjugate poles}
Given $X(z)=\dfrac{1-0.9\cos(0.3\pi)z^{-1}}{1-2(0.9)\cos(0.3\pi)z^{-1}+0.81z^{-2}}$, ROC $|z|>0.9$.

Comparing with $\dfrac{1-r\cos\omega_0\,z^{-1}}{1-2r\cos\omega_0\,z^{-1}+r^{2}z^{-2}}$: $r=0.9,\;\omega_0=0.3\pi,\;r^{2}=0.81.\;\checkmark$ Poles at $z=0.9\,e^{\pm j0.3\pi}$, both with $|z|=0.9<1$.
\begin{equation*}
\boxed{x[n]=(0.9)^{n}\cos(0.3\pi n)\,u[n].}
\end{equation*}
This is a damped sinusoid: $(0.9)^{n}$ envelope decays ($r<1$); oscillation frequency $\omega_0=0.3\pi$ rad/sample.

%=============================================================
\section*{Problem 3: $z$-Transform Properties}
%=============================================================

\subsection*{Part (a): Time shifting}
Given $(0.8)^{n}u[n]\zt\dfrac{1}{1-0.8z^{-1}}$, $|z|>0.8$. By $x[n-2]\zt z^{-2}X(z)$:
\begin{equation*}
\boxed{Y(z)=\dfrac{z^{-2}}{1-0.8\,z^{-1}}, \qquad |z|>0.8.}
\end{equation*}
The factor $z^{-2}$ contributes \textbf{two zeros at $z=0$} (in the $z^{-1}$ form). Rewriting $Y(z)=\dfrac{1}{z^{2}-0.8z}=\dfrac{1}{z(z-0.8)}$ shows one pole at $z=0.8$ and one pole at $z=0$, with the other zero at $z=0$ cancelling one of the shifted poles; the remaining ``missing'' zero sits at $z=\infty$.

\subsection*{Part (b): $z$-domain scaling}
Given $u[n]\zt\dfrac{1}{1-z^{-1}}$, $|z|>1$. Scaling: $z_0^{n}x[n]\zt X(z/z_0)$. With $z_0=0.5$:
\[
(0.5)^{n}u[n]\zt \frac{1}{1-(z/0.5)^{-1}}=\frac{1}{1-0.5z^{-1}}, \qquad |z|>0.5.
\]
\begin{equation*}
\boxed{(0.5)^{n}u[n]\zt\dfrac{1}{1-0.5\,z^{-1}}, \qquad |z|>0.5.}
\end{equation*}
This matches the standard pair with $a=0.5.\;\checkmark$

\subsection*{Part (c): Convolution property}
$h_1[n]=(0.3)^{n}u[n]$, $h_2[n]=(0.7)^{n}u[n]$.

\textbf{(i)} $H_1(z)=\dfrac{1}{1-0.3z^{-1}}$, $|z|>0.3$; $H_2(z)=\dfrac{1}{1-0.7z^{-1}}$, $|z|>0.7$.

\textbf{(ii)} $H(z)=\dfrac{1}{(1-0.3z^{-1})(1-0.7z^{-1})}$, $|z|>0.7$.

\textbf{(iii) Partial fractions:}
\begin{itemize}[leftmargin=*]
\item $z^{-1}=10/3$: $1=A(-4/3)\Rightarrow A=-3/4$.
\item $z^{-1}=10/7$: $1=B(4/7)\Rightarrow B=7/4$.
\end{itemize}
\begin{equation*}
\boxed{h[n]=\left[-\dfrac{3}{4}(0.3)^{n}+\dfrac{7}{4}(0.7)^{n}\right]u[n].}
\end{equation*}

\textbf{(iv) Direct-convolution check:}
\begin{itemize}[leftmargin=*]
\item $h[0]=h_1[0]h_2[0]=1$. Formula: $-3/4+7/4=1.\;\checkmark$
\item $h[1]=h_1[0]h_2[1]+h_1[1]h_2[0]=0.7+0.3=1.0$. Formula: $-0.225+1.225=1.0.\;\checkmark$
\end{itemize}

\subsection*{Part (d): Initial-value theorem and full inversion}
Given $X(z)=\dfrac{3-z^{-1}}{(1-0.4z^{-1})(1+0.5z^{-1})}$, $|z|>0.5$.

\textbf{(i) Initial value:}
\begin{equation*}
\boxed{x[0]=\lim_{z\to\infty}X(z)=\dfrac{3-0}{1\cdot 1}=3.}
\end{equation*}

\textbf{(ii) Partial fractions.} Poles at $z=0.4,\;z=-0.5$.
\begin{itemize}[leftmargin=*]
\item $z^{-1}=2.5$: $3-2.5=A(1+1.25)=2.25A\Rightarrow A=2/9$.
\item $z^{-1}=-2$: $3+2=B(1+0.8)=1.8B\Rightarrow B=25/9$.
\end{itemize}
\[
x[n]=\tfrac{2}{9}(0.4)^{n}u[n]+\tfrac{25}{9}(-0.5)^{n}u[n].
\]
Check: $x[0]=2/9+25/9=3.\;\checkmark$

\textbf{(iii)} $x[1]=\tfrac{2}{9}(0.4)+\tfrac{25}{9}(-0.5)=\tfrac{0.8-12.5}{9}=-\tfrac{11.7}{9}=-1.3$.

$x[2]=\tfrac{2}{9}(0.16)+\tfrac{25}{9}(0.25)=\tfrac{0.32+6.25}{9}\approx 0.73$.

%=============================================================
\section*{Problem 4: System Analysis with $H(z)$}
%=============================================================

\noindent System: $y[n]-0.9\,y[n-1]+0.18\,y[n-2]=x[n]$.

\subsection*{Part (a): System function}
$(1-0.9z^{-1}+0.18z^{-2})Y(z)=X(z)$. Factoring $z^{2}-0.9z+0.18=0\Rightarrow z=\dfrac{0.9\pm 0.3}{2}$, poles at $z=0.6,\;z=0.3$. No finite zeros.
\begin{equation*}
\boxed{H(z)=\dfrac{1}{(1-0.6z^{-1})(1-0.3z^{-1})}, \qquad |z|>0.6 \text{ (causal)}.}
\end{equation*}

\subsection*{Part (b): Stability}
Both poles satisfy $|0.6|<1$ and $|0.3|<1$. Causal system with all poles strictly inside the unit circle $\Rightarrow$ \textbf{BIBO stable}.$\;\checkmark$

\subsection*{Part (c): Impulse response}
Partial fractions: $\dfrac{1}{(1-0.6z^{-1})(1-0.3z^{-1})}=\dfrac{A}{1-0.6z^{-1}}+\dfrac{B}{1-0.3z^{-1}}$.
\begin{itemize}[leftmargin=*]
\item $z^{-1}=5/3$: $1=A(1-0.5)=0.5A\Rightarrow A=2$.
\item $z^{-1}=10/3$: $1=B(1-2)=-B\Rightarrow B=-1$.
\end{itemize}
\begin{equation*}
\boxed{h[n]=\bigl[2(0.6)^{n}-(0.3)^{n}\bigr]u[n].}
\end{equation*}
Check: $h[0]=2-1=1$, matching the coefficient of $x[n]$ in the ODE.$\;\checkmark$

\subsection*{Part (d): Output for input $(0.5)^{n}u[n]$}
\textbf{(i)} $X(z)=\dfrac{1}{1-0.5z^{-1}}$, $|z|>0.5$.
\[
Y(z)=\frac{1}{(1-0.6z^{-1})(1-0.3z^{-1})(1-0.5z^{-1})}, \qquad |z|>0.6.
\]
\textbf{(ii) Partial fractions} $\dfrac{A}{1-0.6z^{-1}}+\dfrac{B}{1-0.3z^{-1}}+\dfrac{C}{1-0.5z^{-1}}$:
\begin{itemize}[leftmargin=*]
\item $z^{-1}=5/3$: $1=A(0.5)(1/6)=A/12\Rightarrow A=12$.
\item $z^{-1}=10/3$: $1=B(-1)(-2/3)=2B/3\Rightarrow B=3/2$.
\item $z^{-1}=2$: $1=C(-0.2)(0.4)=-0.08C\Rightarrow C=-12.5$.
\end{itemize}
Consistency: $A+B+C=12+1.5-12.5=1=Y(\infty)=y[0]$, matching $y[0]=x[0]=1$ from the ODE with zero ICs.$\;\checkmark$
\begin{equation*}
\boxed{y[n]=\bigl[12\,(0.6)^{n}+1.5\,(0.3)^{n}-12.5\,(0.5)^{n}\bigr]u[n].}
\end{equation*}

\textbf{(iii) Verify:}
\begin{itemize}[leftmargin=*]
\item $y[0]=12+1.5-12.5=1$. ODE: $y[0]=x[0]=1.\;\checkmark$
\item $y[1]=12(0.6)+1.5(0.3)-12.5(0.5)=7.2+0.45-6.25=1.4$. ODE: $y[1]=0.9(1)+0.5=1.4.\;\checkmark$
\end{itemize}

\subsection*{Part (e): Stability of $G(z)=\dfrac{1}{1-1.5z^{-1}}$}
Pole at $z=1.5$, outside the unit circle ($|1.5|>1$).

\textbf{If causal:} ROC $|z|>1.5$. Unit circle \emph{not} in ROC. \textbf{Unstable}; $g[n]=(1.5)^{n}u[n]$ diverges; DTFT does not exist.

\textbf{If anti-causal:} ROC $|z|<1.5$. Unit circle \emph{is} in ROC. \textbf{BIBO stable}; $g[n]=-(1.5)^{n}u[-n-1]$ decays as $n\to-\infty$; DTFT exists.

\medskip\noindent\textbf{Key Point.} This parallels the Laplace result from Problem 4(e) of the Lectures 16--18 homework: a pole outside the unit circle makes a causal system unstable but is compatible with a stable anti-causal system. The ROC determines stability, not the pole location alone.

%=============================================================
\section*{Problem 5: The Unilateral $z$-Transform}
%=============================================================

\subsection*{Part (a): First-order with IC}
System: $y[n]-0.8\,y[n-1]=(0.5)^{n}u[n]$, $y[-1]=3$.

\textbf{(i) Unilateral $z$-transform.} Using $y[n-1]\ztu z^{-1}Y(z)+y[-1]$:
\[
\text{LHS: } Y(z)-0.8\bigl[z^{-1}Y(z)+3\bigr]=(1-0.8z^{-1})Y(z)-2.4. \qquad \text{RHS: } \tfrac{1}{1-0.5z^{-1}}.
\]

\textbf{(ii) Solve for $Y(z)$:}
\[
(1-0.8z^{-1})Y(z)=\frac{1}{1-0.5z^{-1}}+2.4\;\Rightarrow\; Y(z)=\frac{1}{(1-0.5z^{-1})(1-0.8z^{-1})}+\frac{2.4}{1-0.8z^{-1}}.
\]

\textbf{(iii) Partial fractions on the first term:}
\begin{itemize}[leftmargin=*]
\item $z^{-1}=2$: $1=A(1-1.6)=-0.6A\Rightarrow A=-5/3$.
\item $z^{-1}=1.25$: $1=B(1-0.625)=0.375B\Rightarrow B=8/3$.
\end{itemize}
So
\[
Y(z)=\frac{-5/3}{1-0.5z^{-1}}+\frac{8/3+12/5}{1-0.8z^{-1}}, \qquad \tfrac{8}{3}+\tfrac{12}{5}=\tfrac{40+36}{15}=\tfrac{76}{15}.
\]
\begin{equation*}
\boxed{y[n]=\left[-\dfrac{5}{3}(0.5)^{n}+\dfrac{76}{15}(0.8)^{n}\right]u[n].}
\end{equation*}

\textbf{(iv) Verify:}
\begin{itemize}[leftmargin=*]
\item $y[0]=-5/3+76/15=-25/15+76/15=51/15=17/5=3.4$. ODE: $y[0]=0.8(3)+1=3.4.\;\checkmark$
\item $y[1]=-5/6+60.8/15\approx-0.833+4.053=3.22$. ODE: $y[1]=0.8(3.4)+0.5=3.22.\;\checkmark$
\end{itemize}

\subsection*{Part (b): Second-order zero-input}
System: $y[n]-0.5\,y[n-1]-0.06\,y[n-2]=0$, $y[-1]=5,\;y[-2]=0$.

\textbf{(i) Unilateral $z$-transforms:}
\[
\mathcal{Z}_u\{y[n-1]\}=z^{-1}Y+5, \qquad \mathcal{Z}_u\{y[n-2]\}=z^{-2}Y+0+5z^{-1}.
\]
Substitute:
\[
Y-0.5(z^{-1}Y+5)-0.06(z^{-2}Y+5z^{-1})=0 \Rightarrow (1-0.5z^{-1}-0.06z^{-2})Y=2.5+0.3z^{-1}.
\]
Factor: $z^{2}-0.5z-0.06=0 \Rightarrow z=\dfrac{0.5\pm 0.7}{2}$, poles at $z=0.6,\;z=-0.1$. Thus
\[
Y(z)=\frac{2.5+0.3z^{-1}}{(1-0.6z^{-1})(1+0.1z^{-1})}.
\]

\textbf{Partial fractions:}
\begin{itemize}[leftmargin=*]
\item $z^{-1}=5/3$: $2.5+0.5=A(1+1/6)=\tfrac{7}{6}A\Rightarrow A=18/7$.
\item $z^{-1}=-10$: $2.5-3=B(1+6)=7B\Rightarrow B=-1/14$.
\end{itemize}
\begin{equation*}
\boxed{y[n]=\left[\dfrac{18}{7}(0.6)^{n}-\dfrac{1}{14}(-0.1)^{n}\right]u[n].}
\end{equation*}
Check: $y[0]=18/7-1/14=36/14-1/14=35/14=2.5$. ODE: $y[0]=0.5(5)+0.06(0)=2.5.\;\checkmark$

\textbf{(iii)} Input is zero $\Rightarrow$ this is a \textbf{pure zero-input response}: output is entirely due to $y[-1]=5$ and $y[-2]=0$.

\subsection*{Part (c): ZSR/ZIR decomposition for part (a)}
From part (a): $Y(z)=\dfrac{1}{(1-0.5z^{-1})(1-0.8z^{-1})}+\dfrac{2.4}{1-0.8z^{-1}}.$

\textbf{Zero-state response} ($y[-1]=0$, keep input): $Y_{\text{ZS}}(z)=\dfrac{1}{(1-0.5z^{-1})(1-0.8z^{-1})}$. Using $A=-5/3,\;B=8/3$ from part (a):
\begin{equation*}
\boxed{y_{\text{ZS}}[n]=\left[-\dfrac{5}{3}(0.5)^{n}+\dfrac{8}{3}(0.8)^{n}\right]u[n].}
\end{equation*}

\textbf{Zero-input response} (input $=0$, keep $y[-1]=3$): $(1-0.8z^{-1})Y_{\text{ZI}}=2.4\Rightarrow Y_{\text{ZI}}(z)=\dfrac{2.4}{1-0.8z^{-1}}$.
\begin{equation*}
\boxed{y_{\text{ZI}}[n]=2.4\,(0.8)^{n}u[n].}
\end{equation*}

\textbf{Verify sum:}
\[
y_{\text{ZS}}+y_{\text{ZI}}=\left[-\tfrac{5}{3}(0.5)^{n}+\bigl(\tfrac{8}{3}+\tfrac{12}{5}\bigr)(0.8)^{n}\right]u[n]=\left[-\tfrac{5}{3}(0.5)^{n}+\tfrac{76}{15}(0.8)^{n}\right]u[n].
\]
This matches $y[n]$ from part (a).$\;\checkmark$

\medskip\noindent\textbf{Key Point.} The ZSR starts at $y_{\text{ZS}}[0]=-5/3+8/3=1=x[0]$ (system at rest). The ZIR starts at $y_{\text{ZI}}[0]=2.4$ (the IC contribution). Their sum gives $y[0]=1+2.4=3.4$, matching the ODE check.

%=============================================================
\section*{Problem Index}
%=============================================================

\begin{itemize}[leftmargin=*]
\item \textbf{Problem 1.} Compute the $z$-transform and ROC for seven signals (causal geometric, anti-causal geometric, delta combinations, $n\cdot a^{n}u[n]$, sum of two geometrics, a two-sided signal, and a damped cosine), then verify $X(1)=\sum x[n]$.
\item \textbf{Problem 2.} Inverse $z$-transforms via partial fractions: distinct real poles causal, distinct real poles two-sided (annular ROC), repeated poles, and complex-conjugate poles (damped cosine).
\item \textbf{Problem 3.} $z$-transform properties --- time shifting, $z$-domain scaling, convolution (cascade), and the initial-value theorem.
\item \textbf{Problem 4.} Second-order causal LTI difference equation: $H(z)$, BIBO stability, $h[n]$, driven output with $(0.5)^{n}u[n]$, and a second system with a pole outside the unit circle under causal vs.\ anti-causal assumptions.
\item \textbf{Problem 5.} Unilateral $z$-transform --- first-order with nonzero IC, second-order zero-input response, and ZSR/ZIR decomposition.
\end{itemize}

\end{document}
